\paragraph{奇商偏序下的零余类压缩生成、唯一极小覆盖与精确计数}
在定理 \ref{thm:xi-time-part71-zero-coset-finite-intersection-pairwise-compatibility} 已给出零余类族任意有限交截的完全分类、定理 \ref{thm:xi-time-part72-zero-coset-divisibility-nested-disjoint-dichotomy} 已把整除链上的交叠压缩为奇倍嵌套与偶倍排斥二分律之后，这条主线还可继续向“表示压缩”推进一步：零点集合 \(\mathcal Z_m\) 并不需要由全部候选余类来生成，而可唯一压缩到奇商偏序下的极大元族；相应的精确计数也可从全约数格容斥进一步压缩为该极大生成族上的包含--排斥闭式。

\begin{theorem}[零余类之间的全局包含关系由奇商整除完全刻画]\label{thm:xi-time-part72a-zero-coset-global-inclusion-odd-divisibility}
\leanpartial{paper\_xi\_time\_part72a\_zero\_coset\_global\_inclusion\_odd\_divisibility}{the requested Lean signature omits the paper's hypothesis that \(M\) is even; at \(M=0\), \(g=1\), and \(h=2\), both zero-cosets equal \(\{0\}\) but \(h/g=2\) is not odd}
设 \(M\) 为偶整数，且 \(g,h\mid M/2\)。则
\[
S_g(M)\subseteq S_h(M)
\iff
g\mid h\ \text{且}\ \frac{h}{g}\ \text{为奇数}.
\]
\end{theorem}
\begin{proof}
若 \(g\mid h\) 且 \(h/g\) 为奇数，则结论已由定理 \ref{thm:xi-time-part72-zero-coset-divisibility-nested-disjoint-dichotomy} 的第一条给出。

反之，设 \(S_g(M)\subseteq S_h(M)\)。则
\[
S_g(M)\cap S_h(M)=S_g(M),
\]
因而交集非空。由定理 \ref{thm:xi-time-part71-zero-coset-finite-intersection-pairwise-compatibility} 在 \(r=2\) 的情形，
\[
\left|S_g(M)\cap S_h(M)\right|=\gcd(g,h).
\]
又由定理 \ref{thm:fold-zero-coset-rigidity}，
\[
|S_g(M)|=g.
\]
于是
\[
\gcd(g,h)=|S_g(M)\cap S_h(M)|=|S_g(M)|=g,
\]
从而 \(g\mid h\)。再由交集非空判据，
\[
2\gcd(g,h)\mid (h-g),
\]
即
\[
2g\mid g\left(\frac{h}{g}-1\right).
\]
故 \(2\mid (h/g-1)\)，也就是 \(h/g\) 为奇数。证毕。
\end{proof}

\begin{theorem}[谱零点集合在奇商偏序下的唯一极小生成族]\label{thm:xi-time-part72a-zero-spectrum-maximal-odd-divisibility-compression}
\leanverified{paper\_xi\_time\_part72a\_zero\_spectrum\_maximal\_odd\_divisibility\_compression}
沿用推论 \ref{cor:fold-zero-coset-union} 与推论 \ref{cor:xi-fold-zero-set-v2-semilattice-decomposition} 的记号。记
\[
M:=F_{m+2},
\qquad
\mathcal D_m:=
\left\{
d:\ d\mid (m+2),\ d<m+2,\ F_d\mid \frac{M}{2}
\right\},
\]
\[
\mathcal G_m:=\{F_d:\ d\in \mathcal D_m\}.
\]
则
\[
\mathcal Z_m=\bigcup_{g\in \mathcal G_m}S_g(M).
\]
在 \(\mathcal G_m\) 上定义偏序
\[
g\preceq h
\iff
S_g(M)\subseteq S_h(M),
\]
等价地，由定理 \ref{thm:xi-time-part72a-zero-coset-global-inclusion-odd-divisibility}，
\[
g\preceq h
\iff
g\mid h\ \text{且}\ \frac{h}{g}\ \text{为奇数}.
\]
记 \(\mathcal M_m\subseteq \mathcal G_m\) 为全部 \(\preceq\)-极大元组成的集合。则成立：
\begin{enumerate}
  \item 谱零点集合具有压缩表示
  \[
  \mathcal Z_m=\bigcup_{g\in \mathcal M_m}S_g(M).
  \]
  \item \(\mathcal M_m\) 是满足上式的唯一极小生成族：若 \(\mathcal A\subseteq \mathcal G_m\) 也满足
  \[
  \mathcal Z_m=\bigcup_{g\in \mathcal A}S_g(M),
  \]
  则必有
  \[
  \mathcal M_m\subseteq \mathcal A.
  \]
\end{enumerate}
\end{theorem}
\begin{proof}
第一条只需删除全部非极大元。事实上，若 \(g\in \mathcal G_m\setminus \mathcal M_m\)，则存在 \(h\in \mathcal G_m\) 使 \(g\prec h\)。由定理 \ref{thm:xi-time-part72a-zero-coset-global-inclusion-odd-divisibility}，
\[
S_g(M)\subseteq S_h(M),
\]
故从并集中去掉 \(S_g(M)\) 不改变并集。对所有非极大元依次执行此操作，即得
\[
\mathcal Z_m=\bigcup_{g\in \mathcal M_m}S_g(M).
\]

现证唯一极小性。任取 \(g\in \mathcal M_m\)。若 \(h\in \mathcal G_m\setminus\{g\}\) 且
\[
S_g(M)\cap S_h(M)\neq \varnothing,
\]
则由定理 \ref{thm:xi-time-part71-zero-coset-finite-intersection-pairwise-compatibility}，
\[
S_g(M)\cap S_h(M)=S_{\gcd(g,h)}(M)
\]
且
\[
\gcd(g,h)<g,
\]
否则若 \(\gcd(g,h)=g\)，则 \(g\mid h\)，再由交非空判据可得 \(h/g\) 为奇数，从而 \(g\prec h\)，与 \(g\) 的极大性矛盾。

取 \(S_g(M)\) 的标准代表点
\[
k_g:=\frac{M}{2g}\in S_g(M).
\]
若 \(k_g\in S_h(M)\) 对某个 \(h\neq g\) 成立，则由上式知
\[
k_g\in S_g(M)\cap S_h(M)=S_{\gcd(g,h)}(M)
\]
且 \(d:=\gcd(g,h)<g\)。由于交集非空，\(g/d\) 必为奇数。于是 \(S_d(M)\subseteq S_g(M)\)，并且按定理 \ref{thm:fold-zero-coset-rigidity} 的显式描述，\(S_d(M)\) 在 \(S_g(M)\) 内对应于模 \(g/d\) 的单一余类
\[
t\equiv \frac{g/d-1}{2}\pmod{g/d}.
\]
而点 \(k_g\) 对应于参数 \(t=0\)。因为 \(g/d>1\) 且为奇数，
\[
\frac{g/d-1}{2}\not\equiv 0\pmod{g/d},
\]
故 \(k_g\notin S_d(M)\)，矛盾。

因此
\[
k_g\notin \bigcup_{h\in \mathcal G_m\setminus\{g\}}S_h(M),
\]
从而
\[
S_g(M)\nsubseteq \bigcup_{h\in \mathcal G_m\setminus\{g\}}S_h(M).
\]
于是任何生成整个 \(\mathcal Z_m\) 的子族都必须包含 \(g\)。由于 \(g\in \mathcal M_m\) 任意，遂得
\[
\mathcal M_m\subseteq \mathcal A.
\]
这就证明了唯一极小性。证毕。
\end{proof}

\begin{theorem}[极大奇商生成族上的谱零点集合精确容斥公式]\label{thm:xi-time-part72a-zero-spectrum-maximal-family-inclusion-exclusion}
沿用定理 \ref{thm:xi-time-part72a-zero-spectrum-maximal-odd-divisibility-compression} 的记号，写
\[
\mathcal M_m=\{g_1,\dots,g_r\}.
\]
则谱零点集合的基数满足精确公式
\[
|\mathcal Z_m|
=
\sum_{\varnothing\neq I\subseteq \{1,\dots,r\}}
(-1)^{|I|+1}
\mathbf 1_{\mathrm{comp}(I)}
\gcd(g_i:\ i\in I),
\]
其中相容指标 \(\mathbf 1_{\mathrm{comp}(I)}\) 定义为：当且仅当对全部 \(i,j\in I\) 都有
\[
2\gcd(g_i,g_j)\mid (g_j-g_i)
\]
时，\(\mathbf 1_{\mathrm{comp}(I)}=1\)；否则 \(\mathbf 1_{\mathrm{comp}(I)}=0\)。
\end{theorem}
\begin{proof}
由定理 \ref{thm:xi-time-part72a-zero-spectrum-maximal-odd-divisibility-compression}，
\[
\mathcal Z_m=\bigcup_{\nu=1}^{r}S_{g_\nu}(M).
\]
对该有限并应用包含--排斥，
\[
|\mathcal Z_m|
=
\sum_{\varnothing\neq I\subseteq \{1,\dots,r\}}
(-1)^{|I|+1}
\left|
\bigcap_{i\in I}S_{g_i}(M)
\right|.
\]
再由定理 \ref{thm:xi-time-part71-zero-coset-finite-intersection-pairwise-compatibility}，交截非空当且仅当上述两两相容条件成立，而一旦非空就有
\[
\left|
\bigcap_{i\in I}S_{g_i}(M)
\right|
=
\gcd(g_i:\ i\in I).
\]
代回即得所示公式。证毕。
\end{proof}

\begin{corollary}[极大生成族 pairwise-incompatible 时的精确直和公式]\label{cor:xi-time-part72a-zero-spectrum-maximal-family-pairwise-incompatible-exact-sum}
沿用定理 \ref{thm:xi-time-part72a-zero-spectrum-maximal-odd-divisibility-compression} 的记号。若极大生成族 \(\mathcal M_m=\{g_1,\dots,g_r\}\) 满足任意两元都不相容，即
\[
2\gcd(g_i,g_j)\nmid (g_j-g_i)
\qquad (i\neq j),
\]
则对应余类两两不交，从而
\[
|\mathcal Z_m|=\sum_{\nu=1}^{r}g_\nu.
\]
\end{corollary}
\begin{proof}
由定理 \ref{thm:xi-time-part71-zero-coset-finite-intersection-pairwise-compatibility}，上述条件等价于
\[
S_{g_i}(M)\cap S_{g_j}(M)=\varnothing
\qquad (i\neq j).
\]
故
\[
\mathcal Z_m=\bigsqcup_{\nu=1}^{r}S_{g_\nu}(M).
\]
再由定理 \ref{thm:fold-zero-coset-rigidity} 的 \(|S_g(M)|=g\)，可得
\[
|\mathcal Z_m|
=
\sum_{\nu=1}^{r}|S_{g_\nu}(M)|
=
\sum_{\nu=1}^{r}g_\nu.
\]
证毕。
\end{proof}

\noindent
上述结果表明，零余类族的算术几何还具有一层此前未被完全显化的“压缩表示”刚性：有限交截虽由全体候选余类的 pairwise 相容性控制，但真正生成 \(\mathcal Z_m\) 的并不需要整张约数图，而只需要奇商偏序下的极大元族；这一压缩不仅是存在性的，而且是唯一极小的。相应地，谱零点计数也可从全约数格的指标容斥，进一步压缩为极大生成族上的精确包含--排斥；而在 pairwise-incompatible 的情形下，全部交叠复杂性彻底消失，精确计数退化为一条纯粹的 Fibonacci 权重直和公式。于是，零余类交叠、包含、压缩表示与 exact counting 四者就在同一条 odd-divisibility 主线上获得了闭合。
